% source: Robin Hartshorne, "Algebraic Geometry", GTM 52 (1977)
% pdf_page: 0243
% doc_page: 226 (Chapter III, Section 5, continuation of the proof of Theorem 5.1)
% retrieved: 2026-07-11
% transcribed_by: claude-sonnet-5 (vision, rendered at 200dpi via PyMuPDF)

\DeclareMathOperator{\Spec}{Spec}

% [continuation of the proof of Theorem 5.1, from page-0242.tex]

of the grading by $n$, so that we can sort out the pieces at the end. Note
that all the cohomology groups in question have a natural structure of
$A$-module (2.6.1).

For each $i = 0, \ldots, r$, let $U_i$ be the open set $D_+(x_i)$. Then each
$U_i$ is an open affine subset of $X$, and the $U_i$ cover $X$, so we can
compute the cohomology of $\mathcal{F}$ by using \v{C}ech cohomology for the
covering $\mathfrak{U} = (U_i)$, by (4.5). For any set of indices
$i_0, \ldots, i_p$, the open set $U_{i_0, \ldots, i_p}$ is just
$D_+(x_{i_0} \cdots x_{i_p})$, so by (II, 5.11) we have
\[
  \mathcal{F}(U_{i_0, \ldots, i_p}) \cong S_{x_{i_0} \cdots x_{i_p}},
\]
the localization of $S$ with respect to the element $x_{i_0} \cdots x_{i_p}$.
Furthermore, the grading on $\mathcal{F}$ corresponds to the natural grading
of $S_{x_{i_0} \cdots x_{i_p}}$ under this isomorphism. Thus the \v{C}ech
complex of $\mathcal{F}$ is given by
\[
  C^{\boldsymbol{\cdot}}(\mathfrak{U}, \mathcal{F}) :
    \prod S_{x_{i_0}} \to \prod S_{x_{i_0} x_{i_1}} \to \cdots \to
    S_{x_0 \cdots x_r},
\]
and the modules all have a natural grading compatible with the grading on
$\mathcal{F}$.

Now $H^0(X, \mathcal{F})$ is the kernel of the first map, which is just $S$,
as we have seen earlier (II, 5.13). This proves (a).

Next we consider $H^r(X, \mathcal{F})$. It is the cokernel of the last map
in the \v{C}ech complex, which is
\[
  d^{r-1} : \prod_k S_{x_0 \cdots \hat{x}_k \cdots x_r} \to
    S_{x_0 \cdots x_r}.
\]
We think of $S_{x_0 \cdots x_r}$ as a free $A$-module with basis
$x_0^{l_0} \cdots x_r^{l_r}$, with $l_i \in \mathbf{Z}$. The image of
$d^{r-1}$ is the free submodule generated by those basis elements for which
at least one $l_i \ge 0$. Thus $H^r(X, \mathcal{F})$ is a free $A$-module
with basis consisting of the ``negative'' monomials
\[
  \{ x_0^{l_0} \cdots x_r^{l_r} \mid l_i < 0 \text{ for each } i \}.
\]
Furthermore the grading is given by $\sum l_i$. There is only one such
monomial of degree $-r-1$, namely $x_0^{-1} \cdots x_r^{-1}$, so we see that
$H^r(X, \mathcal{O}_X(-r-1))$ is a free $A$-module of rank $1$. This proves
(c).

To prove (d), first note that if $n < 0$, then $H^0(X, \mathcal{O}_X(n)) = 0$
by (a), and $H^r(X, \mathcal{O}_X(-n-r-1)) = 0$ by what we have just seen,
since in that case $-n-r-1 > -r-1$, and there are no negative monomials of
that degree. So the statement is trivial for $n < 0$. For $n \ge 0$,
$H^0(X, \mathcal{O}_X(n))$ has a basis consisting of the usual monomials of
degree $n$, i.e., $\{ x_0^{m_0} \cdots x_r^{m_r} \mid m_i \ge 0 \text{ and }
\sum m_i = n \}$. The natural pairing with $H^r(X, \mathcal{O}_X(-n-r-1))$
into $H^r(X, \mathcal{O}_X(-r-1))$ is determined by
\[
  (x_0^{m_0} \cdots x_r^{m_r}) \cdot (x_0^{l_0} \cdots x_r^{l_r}) =
    x_0^{m_0+l_0} \cdots x_r^{m_r+l_r},
\]
where $\sum l_i = -n-r-1$, and the object on the right becomes $0$ if any
$m_i + l_i \ge 0$. So it is clear that we have a perfect pairing, under
which $x_0^{-m_0-1} \cdots x_r^{-m_r-1}$ is the dual basis element of
$x_0^{m_0} \cdots x_r^{m_r}$.
\end{proof}
